%%%% --------------------------------------------------
%%%% PAQUETES BÁSICOS
%%%% --------------------------------------------------
\usepackage[utf8]{inputenc}       % Codificación UTF-8 para soportar acentos y caracteres especiales
\usepackage[spanish]{babel}       % Idioma español, adapta títulos automáticos (Figura, Tabla, etc.)
\usepackage{csquotes}             % Mejora el manejo de comillas y compatibilidad con biblatex
\usepackage{xurl}                 % Permite cortar URLs largas automáticamente en varias líneas
\usepackage{microtype}            % Mejora la tipografía (espaciado, justificación, ligaduras)

%%%% --------------------------------------------------
%%%% MATEMÁTICAS
%%%% --------------------------------------------------
\usepackage{amsmath, amssymb, amsfonts}  % Entornos y símbolos matemáticos avanzados

%%%% --------------------------------------------------
%%%% IMÁGENES Y FIGURAS
%%%% --------------------------------------------------
\usepackage{graphicx}             % Permite incluir imágenes (\includegraphics)
\usepackage{caption}              % Configuración de captions en figuras y tablas
\usepackage{float}                % Permite forzar la posición exacta de figuras/tablas con [H]

%%%% --------------------------------------------------
%%%% TABLAS
%%%% --------------------------------------------------
\usepackage{array}               % Mejoras en el manejo de tablas
\usepackage{multirow}            % Permite combinar filas en tablas
\usepackage{multicol}            % Permite crear columnas múltiples de texto
\usepackage{booktabs}            % Tablas con líneas profesionales (\toprule, \midrule, \bottomrule)

% =============================
% PAQUETES PARA CÓDIGO Y PSEUDOCÓDIGO
% =============================
\usepackage{listings}             % Entorno para mostrar código fuente
\usepackage[linesnumbered,ruled,vlined]{algorithm2e} % Entorno para pseudocódigo

%%%% --------------------------------------------------
%%%% HIPERVÍNCULOS Y PDF
%%%% --------------------------------------------------
\usepackage{hyperref}             % Permite enlaces internos y externos

%%%% --------------------------------------------------
%%%% BIBLIOGRAFÍA
%%%% --------------------------------------------------
\usepackage[backend=biber, style=ieee]{biblatex}  % Gestión de bibliografía estilo IEEE

%%%% --------------------------------------------------
%%%% OTROS PAQUETES ÚTILES
%%%% --------------------------------------------------
\usepackage{comment}              % Permite comentar bloques de código grandes
\usepackage{tcolorbox}            % Permite crear recuadros de colores para resaltar texto
\usepackage{xcolor}               % Permite definir colores personalizados
\usepackage{titlesec}             % Personalización de títulos de secciones y subsecciones
\usepackage{setspace}             % Control de interlineado
\usepackage{fancyhdr}             % Personalización de encabezados y pies de página
\usepackage[perpage]{footmisc}   % Notas al pie personalizadas (reinicia por página, opciones)

%%%% --------------------------------------------------
%%%% CONFIGURACIÓN DE MÁRGENES
%%%% --------------------------------------------------
\usepackage{geometry}             % Permite configurar márgenes
